\section{Stakeholders e Impacto}

\subsection{Stakeholders Identificados}

O sistema proposto impacta diretamente diferentes áreas da organização envolvidas no ciclo de vida das integrações com locadoras, desde o onboarding até a operação contínua.

Os principais stakeholders identificados são:

\begin{itemize}
    \item Equipe RCM e Global
    \item Squad Partners (Engenharia de Integrações)
    \item Equipe de Sustentação (N1)
    \item Liderança Técnica e Arquitetura
\end{itemize}

\subsection{Impacto por Stakeholder}

\subsubsection{RCM e Global}

A equipe RCM e Global é responsável pelo onboarding das locadoras no ecossistema, incluindo cadastro, configuração inicial e preparação para integração técnica, além do suporte pós-implementação.

Atualmente, o processo de integração e manutenção de locadoras pode ser impactado pela complexidade técnica e pelo tempo necessário para validações e ajustes durante a implementação.

Com a solução proposta, espera-se:

\begin{itemize}
    \item Redução no tempo necessário para disponibilização de novas locadoras;
    \item Maior previsibilidade no processo de integração;
    \item Diminuição da dependência de validações manuais e retrabalho técnico;
    \item Aceleração do ciclo entre onboarding e disponibilização de ofertas ao usuário final.
\end{itemize}

Como consequência, a equipe passa a operar com maior eficiência e menor dependência de ciclos longos de validação técnica.

\subsubsection{Squad Partners (Engenharia de Integrações)}

A Squad Partners é responsável pelo desenvolvimento e manutenção das integrações com locadoras, incluindo implementação de novos drivers e fornecedores, correção de bugs e garantia da disponibilidade das ofertas.

Atualmente, a equipe enfrenta desafios relacionados à complexidade das integrações, dificuldade de validação estruturada e fragmentação das ferramentas de observabilidade, o que reduz a visibilidade sobre o comportamento real do sistema.

A solução proposta impacta diretamente a squad, proporcionando:

\begin{itemize}
    \item Maior velocidade no processo de migração de integrações;
    \item Redução do esforço necessário para manutenção de integrações existentes;
    \item Melhoria significativa na capacidade de diagnóstico e investigação de problemas;
    \item Centralização de informações de observabilidade, substituindo múltiplas ferramentas com baixa efetividade;
    \item Introdução de uma camada de visibilidade orientada a negócio, permitindo correlacionar falhas técnicas com impacto real.
\end{itemize}

Esse conjunto de melhorias contribui para aumentar a eficiência da squad, reduzir o retrabalho e elevar a confiabilidade das entregas.

\subsubsection{Equipe de Sustentação (N1)}

A equipe de sustentação é responsável pelo atendimento de primeiro nível, atuando na identificação e resolução inicial de incidentes relacionados ao sistema.

Atualmente, a ausência de visibilidade estruturada sobre o comportamento das integrações dificulta a identificação de problemas e aumenta o tempo de resolução.

Com a solução proposta, a equipe passa a contar com:

\begin{itemize}
    \item Acesso a informações mais claras e centralizadas sobre falhas de integração;
    \item Maior capacidade de diagnóstico inicial sem dependência direta da engenharia;
    \item Redução do tempo médio de resolução de incidentes;
    \item Geração de insumos mais qualificados para escalonamento de problemas.
\end{itemize}

Isso permite uma atuação mais eficiente e autônoma da equipe de sustentação.

\subsubsection{Liderança Técnica e Arquitetura}

A migração para uma arquitetura mais modular é um movimento estratégico de longo prazo, que exige controle de riscos e governança técnica.

A solução proposta oferece:

\begin{itemize}
    \item Maior controle e visibilidade sobre o processo de migração;
    \item Redução de riscos associados à modernização da arquitetura;
    \item Base estruturada para evolução contínua das integrações;
    \item Apoio à tomada de decisão técnica baseada em dados.
\end{itemize}